%!TEX root = ../kr.tex
\chapter[(Довідковий) Словник спеціалізованих термінів]{Словник спеціалізованих термінів}
\label{apdx:text}

Додатки створюються так само, як і основний текст пояснювальної записки.

% Тимчасово прибираємо картинку, якої немає в проекті
%\begin{figure}[!ht]
% \centering
% % \includegraphics{img/Tux.png} 
% \caption{Пінгвін}
% \label{apdxfig:tux}
%\end{figure}

\begin{table}[!ht]
% \underonespace % <-- Ця команда викликала помилку, видаляємо її
\caption{\label{apdxtable:2}Основні та деякі похідні одиниці системи SI}
% Для роботи стовпчиків типу 'm' потрібен пакет array (ми додали його в kr.tex)
\begin{tabular}{|>{\raggedright\arraybackslash}p{5cm}|c|>{\raggedright\arraybackslash}p{4cm}|c|}
\hline
\multicolumn{2}{|c|}{Величина}&\multicolumn{2}{c|}{Одиниця}\\[0pt]\hline
Найменування & Розмірність & Найменування & Позначення \\[0pt]\hline\hline
Довжина & $L$ & метр & м \\[0pt]\hline
Маса & $M$ & кілограм & кг \\[0pt]\hline
Час & $T$ & секунда & с \\[0pt]\hline
\end{tabular}
\end{table}
